\begin{abstract}
An AI agent alternates between model calls and tool calls. While a tool runs, the serving engine keeps the agent's
key--value (KV) cache only if nothing else needs the memory; when it is evicted, the next call recomputes the whole
context. We measure what this costs in GPU energy on an H200 with vLLM and SGLang, three Qwen3 models (32B, 30B-A3B
and 8B) and four workloads: synthetic sessions, recorded SWE-agent trajectories, replayed samples of production
tool-agent traffic from Kimi, and graded BFCL multi-turn tasks. Every prediction was registered before its run.

We report four results. First, on synthetic sessions, once the agents' working set exceeds the cache by a factor of 1.2--2.2, hits collapse;
past that point energy per turn is within 0.01--1.6\% of running without a cache. With
Qwen3-32B and 5\,s pauses the collapse comes between 16 and 32 sessions of about 11k tokens and between 8 and 16 of
about 27k. Second, restoring a prompt's KV from host memory takes about 1/110 to 1/220 of the above-idle GPU energy of
recomputing it. Third, on a 1/4 sample of the production trace, which overloads the GPU, a 160\,GB host tier let the GPU
finish every request (85--86\% with the GPU cache) and process 16--18\% more prompt tokens for about the same energy, in two
time windows; on a lighter 1/8 sample the tier saved under 1\%. Given measured hit rates, a one-parameter model predicts the busy card's energy per
token to within about 1\%. Fourth, on synthetic sessions, an admission gate, switched on only when the working set
exceeds 1.2 times the cache, cost no energy below that point; above it, in two 10-minute runs, it served every session,
completed 48--62\% more turns and cut energy per turn by 42--50\%, at the price of minutes of queueing at
admission. On BFCL, accuracy with and without the cache differed by 3 of 400 tasks, and the
cache doubled correct answers per kWh. All energy is GPU board energy.
\end{abstract}
